% =====================================================================
% MNS4A certificate - all text placed over the background image
%
% Put the background image (logos, MNS4A header, frame; WITHOUT any of
% the text below) in the same folder as fundo.jpg (or change \fundo).
% Compile with: pdflatex certificado_mns4a_texto.tex
%
% Usage (one page per certificate):
%   \certificado{Name}{Paragraph text}
% Use \\ inside the paragraph for the line break.
% The fixed texts (title, "This is to certify that", signatories) are
% defined in the configuration block below.
% =====================================================================
\documentclass[a4paper,landscape]{article}
\usepackage[utf8]{inputenc}
\usepackage[T1]{fontenc}
\usepackage[a4paper,landscape,margin=0cm]{geometry}
\usepackage{graphicx}
\usepackage{xcolor}
\usepackage{tikz}
\usepackage{helvet}
\usepackage{mathpazo}   % serif font for the name
\renewcommand{\familydefault}{\sfdefault}
\pagestyle{empty}
\setlength{\parindent}{0pt}

\definecolor{navy}{HTML}{0B2A5B}
\definecolor{bluetext}{HTML}{1F4E8C}

% ---------- Configuration ----------
\newcommand{\fundo}{fundo.jpg}   % background image

% Fixed texts
\newcommand{\titulo}{CERTIFICATE OF RECOGNITION}
\newcommand{\certifythat}{This is to certify that}

% Signatories
\newcommand{\assinanteA}{Dr. Esfhan Alam Kherani}
\newcommand{\cargoA}{Chair of NMS4A}
\newcommand{\assinanteB}{Dr. L\'aysa Cristina Araujo Resende Chagas}
\newcommand{\cargoB}{President of SBGEA}

% Vertical positions (distance from the top of the page) - adjust if needed
\newcommand{\ytitulo}{8.0cm}
\newcommand{\ycertify}{9.6cm}
\newcommand{\ynome}{11.6cm}
\newcommand{\ycorpo}{14.2cm}
\newcommand{\yassin}{18.0cm}

% ---------- One certificate ----------
% #1 name | #2 paragraph text
\newcommand{\certificado}[2]{%
  \thispagestyle{empty}%
  \begin{tikzpicture}[remember picture,overlay,shift={(current page.north west)}]
    \node[anchor=north west,inner sep=0] at (0,0)
      {\includegraphics[width=\paperwidth,height=\paperheight]{\fundo}};
    % Title
    \node[font=\LARGE\color{bluetext}]
      at (0.5\paperwidth,-\ytitulo) {\titulo};
    % "This is to certify that"
    \node[font=\large\color{navy}]
      at (0.5\paperwidth,-\ycertify) {\certifythat};
    % Name
    \node[font=\fontsize{28}{32}\selectfont\bfseries\rmfamily\color{navy}]
      at (0.5\paperwidth,-\ynome) {#1};
    % Body
    \node[font=\Large\color{navy},text width=0.95\paperwidth,align=center]
      at (0.5\paperwidth,-\ycorpo) {#2};
    % Signatories
    \node[align=center,font=\normalsize\color{navy}]
      at (0.328\paperwidth,-\yassin) {{\bfseries\assinanteA}\\\cargoA};
    \node[align=center,font=\normalsize\color{navy}]
      at (0.715\paperwidth,-\yassin) {{\bfseries\assinanteB}\\\cargoB};
  \end{tikzpicture}%
  \clearpage}

% ---------- Default texts for presentation certificates ----------
% #1 presenter name | #2 title of the work | #3 session name
\newcommand{\localdata}{of the MNS4A Workshop,
  held in S\~ao Jos\'e dos Campos (SP), Brazil, from 19 to 21 October 2026.}

% Oral presentation
\newcommand{\oral}[3]{%
  \certificado{#1}{has presented the oral talk entitled {\bfseries ``#2''}
    in the Session ``#3'' \localdata}}

% Poster presentation
\newcommand{\poster}[3]{%
  \certificado{#1}{has presented the poster entitled {\bfseries ``#2''}
    in the Session ``#3'' \localdata}}

% Invited talk
\newcommand{\convidado}[3]{%
  \certificado{#1}{has presented the invited talk entitled {\bfseries ``#2''}
    in the Session ``#3'' \localdata}}

% ---------- Default text for attendee (listener) certificates ----------
% #1 attendee name
\newcommand{\ouvinte}[1]{{%
  \renewcommand{\titulo}{CERTIFICATE OF PARTICIPATION}%
  \certificado{#1}{has attended the MNS4A Workshop,
    held in S\~ao Jos\'e dos Campos (SP), Brazil, from 19 to 21 October 2026.}}}

\begin{document}

% row 2
\convidado
  {Cassiano Antonio Bortolozo}
  {Geofísica Aplicada A Estudos Ambientais E Desastres Naturais}
  {Modeling of Atmosphere and Ionosphere Variabilities}

% row 3
\convidado
  {Lucie Rolland}
  {Modeling earthquake- and tsunami-induced ionospheric disturbances: from physical mechanisms to source characterization}
  {Modeling of Atmosphere and Ionosphere Variabilities}

% row 4
\convidado
  {Clezio Marcos Denardini}
  {The Instruments for the Rise Campaign for Ionospheric Study}
  {Modeling of Atmosphere and Ionosphere Variabilities}

% row 5
\oral
  {Esfhan Alam Kherani}
  {IONO-DIET with GO-EUREKA: A step towards early warning of catastrophic geohazards}
  {Modeling of Atmosphere and Ionosphere Variabilities}

% row 6
\oral
  {Jonas Rodrigues de Souza}
  {SUPIM-INPE Model: Some Ionospheric Predictions for Equatorial and Low Latitudes}
  {Modeling of Atmosphere and Ionosphere Variabilities}

% row 7
\oral
  {Luis Eduardo Antunes Vieira}
  {Physics-Informed Machine Learning for Radiation Belt Dynamics}
  {Modeling of Atmosphere and Ionosphere Variabilities}

% row 8
\oral
  {Oluwasegun Micheal Adebayo}
  {Earthquake Source Processes Revealed by Dual-Mode Co-Seismic Ionospheric Disturbances}
  {Modeling of Atmosphere and Ionosphere Variabilities}

% row 9
\oral
  {Jose Paulo Marchezi}
  {Can empirical wave models resolve what drives energetic electron precipitation? A comparison of model-driven and measurement-driven diffusion coefficients}
  {Modeling of Atmosphere and Ionosphere Variabilities}

% row 10
\oral
  {Rogério Hisashi Honda}
  {A Thunderstorm--Atmosphere--Ionosphere Coupling Model for the Ionospheric Signatures of Acoustic-Gravity Waves}
  {Modeling of Atmosphere and Ionosphere Variabilities}

% row 11
\oral
  {Sophia Rodrigues Laranja}
  {A Data-Driven Approach to Improving the Representation of Equatorial Vertical Drifts over South America: Preliminary Results on Atmospheric Tides}
  {Modeling of Atmosphere and Ionosphere Variabilities}

% row 12
\oral
  {André Ricardo Fasanaro Martinon}
  {Real-Time Maps of S4 and $\sigma$$\varphi$ Scintillation Indices Over Brazil}
  {Modeling of Atmosphere and Ionosphere Variabilities}

% row 13
\poster
  {Giorgio Picanço}
  {Modeling the Topside Ionosphere Using an Ionosonde-Based Assimilative Model}
  {Modeling of Atmosphere and Ionosphere Variabilities}

% row 14
\poster
  {Thiago Moeda Sant Anna}
  {Modeling the Variability of Relativistic Electron Flux in the Outer Radiation Belt Using LSTM Recurrent Neural Network}
  {Modeling of Atmosphere and Ionosphere Variabilities}

% row 15
\poster
  {Akinbode Oluwasola Olakitan}
  {Machine Learning Applied For Prediction of Maximum Usable Frequency (MUF) in Brazilian Ionospheric Stations}
  {Modeling of Atmosphere and Ionosphere Variabilities}

% row 16
\poster
  {Atotuomah Eyitope Solomon}
  {Updating the Studies of Cosmic Noise Absorption Across the SAMA Based on Recent Measurements}
  {Modeling of Atmosphere and Ionosphere Variabilities}

% row 17
\poster
  {Diego Barros Silva}
  {Asymmetric Evolution of Equatorial Plasma Bubbles at Geomagnetically Conjugate Sites in Brazil}
  {Modeling of Atmosphere and Ionosphere Variabilities}

% row 18
\poster
  {Gowher Rather}
  {Lead-Time-Dependent Machine Learning Prediction of Equatorial Plasma Bubble Onset over Brazil: Snapshot, Temporal, and Spatial Model}
  {Modeling of Atmosphere and Ionosphere Variabilities}

% row 19
\poster
  {Laura Luiz Trigo}
  {Development of Computational Tools in MATLAB for GNSS-Based Ionospheric Analysis Using PCA and Time-Distance Diagrams}
  {Modeling of Atmosphere and Ionosphere Variabilities}

% row 20
\poster
  {Laysa Resende}
  {MUF Index Variability and Its Impacts over Brazilian Stations During a Intense Geomagnetic Storm}
  {Modeling of Atmosphere and Ionosphere Variabilities}

% row 21
\convidado
  {Elvira Astafyeva}
  {GO-EUREKA: toward near-real-time detection and assessment of natural hazards from the ionosphere}
  {Scientific Networks for Atmosphere and Ionospheric Applications}

% row 22
\convidado
  {Maria Graciela Molina}
  {Machine Learning and Ionospheric Datasets: From Data Preparation to Operational Forecasting}
  {Scientific Networks for Atmosphere and Ionospheric Applications}

% row 23
\convidado
  {Michela Ravanelli}
  {Multi-GNSS Networks for Real-Time Monitoring of Atmosphere--Ionosphere Coupling and Natural Hazards}
  {Scientific Networks for Atmosphere and Ionospheric Applications}

% row 24
\oral
  {Alexandre Alvares Pimenta}
  {Bolide Explosions and Climate Change in the Upper Atmosphere: Possible Implications for VLEO/LEO Satellites and Space Debris}
  {Scientific Networks for Atmosphere and Ionospheric Applications}

% row 25
\oral
  {Angela Santos}
  {The Brazilian Digisonde and Fabry--Perot Network: Current Status, Scientific Applications, and Perspectives for Expansion}
  {Scientific Networks for Atmosphere and Ionospheric Applications}

% row 26
\oral
  {Carolina de Sousa do Carmo}
  {Unexpected Suppression and Longitudinal Asymmetry of Equatorial Plasma Bubbles during the 7--8 October 2024 Geomagnetic Storm over the American Sector}
  {Scientific Networks for Atmosphere and Ionospheric Applications}

% row 27
\oral
  {Danny Scipion}
  {Multi-Instrument Ground Networks for Equatorial Aeronomy}
  {Scientific Networks for Atmosphere and Ionospheric Applications}

% row 28
\oral
  {Gabriel Brando Soares}
  {Modernization, Structure and Applications of the Observatório Nacional Geomagnetic Network}
  {Scientific Networks for Atmosphere and Ionospheric Applications}

% row 29
\oral
  {Olufemi Louis Ogunmola}
  {Cusp-Type Sporadic-E Layers at São José dos Campos and Palmas Low-Latitude Stations: Morphology, Altitude Distribution and Tidal Modulation}
  {Scientific Networks for Atmosphere and Ionospheric Applications}

% row 30
\oral
  {Saul A. Sanchez}
  {NRT-TECPRO: Near-Real-Time GNSS Ionospheric Monitoring for Natural Hazard Detection}
  {Scientific Networks for Atmosphere and Ionospheric Applications}

% row 31
\oral
  {Vania Fatima Andrioli Novaes Silva}
  {Monitoring Atmospheric Responses to Natural Hazards Using Meteor Radar: Detection of the Lamb Wave from the 2022 Hunga Tonga Eruption}
  {Scientific Networks for Atmosphere and Ionospheric Applications}

% row 32
\poster
  {Anita Jungton Valcorte}
  {Statistical Analysis of GNSS-Derived Ionospheric Disturbance Indices Over the SAMA Region}
  {Scientific Networks for Atmosphere and Ionospheric Applications}

% row 33
\poster
  {Guillermo Caro}
  {Preliminary results of the study of the long-term effects on the diurnal variation of the geomagnetic field in the SAMA region}
  {Scientific Networks for Atmosphere and Ionospheric Applications}

% row 34
\poster
  {Oladayo Olayiwola Afolabi}
  {From the Equator to Midlatitudes: The Magnetic Footprint of the May 2024 Disturbance Dynamo over South America}
  {Scientific Networks for Atmosphere and Ionospheric Applications}

% row 35
\poster
  {Virginia Klausner}
  {Seismo-Ionospheric Response to the August 22, 2025 Mw 7.5 Drake Passage Earthquake: Spatial Patterns and Propagation Characteristics}
  {Scientific Networks for Atmosphere and Ionospheric Applications}

% row 36
\poster
  {Gbadebo Seun Adeniran}
  {Severe Scintillation Intensification and Large-to-Mesoscale Wave Structures During the 2023 Major Sudden Stratospheric Warming}
  {Scientific Networks for Atmosphere and Ionospheric Applications}

% row 37
\convidado
  {Boris Maletckii}
  {Understanding Seismo-Atmosphere-Ionosphere Coupling by Analysis of Spatio-Temporal Characteristics of Ionospheric Disturbances}
  {Scientific Instrumentation and Challenges}

% row 38
\convidado
  {George Sand L. A. de França}
  {Distributed Acoustic Sensing (DAS) in Seismology: Advances Over the Past Decade, Offshore Applications, and the Brazilian Perspective.}
  {Scientific Instrumentation and Challenges}

% row 39
\convidado
  {Joao Francisco Galera Monico}
  {INCT GNSS NavAer: History, Contributions, and Applications in Space Weather}
  {Scientific Instrumentation and Challenges}

% row 40
\oral
  {Rodolfo de Jesus}
  {Electrodynamics and ionospheric irregularities during a weak geomagnetic storm of 15 April 2025 and intense storm of 16 April 2025 over Brazil}
  {Scientific Instrumentation and Challenges}

% row 41
\oral
  {Danny Scipión}
  {New Jicamarca Radar Capabilities: The All-ISR and JULIA MP modes}
  {Scientific Instrumentation and Challenges}

% row 42
\oral
  {Francisco Carlos de Meneses Junior}
  {Conceptualizing an App for Tracking Aeronomic Conditions over Brazil for Precision Agriculture End-Users}
  {Scientific Instrumentation and Challenges}

% row 43
\oral
  {Cristiano Max Wrasse}
  {Ground-Based Observations of the Mesosphere and Lower Thermosphere: An Overview of the EMBRACE Monitoring Network}
  {Scientific Instrumentation and Challenges}

% row 44
\oral
  {Juan C. Espinoza}
  {J-ARGUS: Expanding Jicamarca's Radar Capabilities}
  {Scientific Instrumentation and Challenges}

% row 45
\oral
  {Rosemary Aparecida Odorizi Lima}
  {Development of Low-Cost Meteorological Instrumentation for Science Education and Public Engagement: The INPE Education Program.}
  {Scientific Instrumentation and Challenges}

% row 46
\oral
  {Sávio Jenner Vianna Silva Rennó}
  {Temporal Climatology of Ionospheric Scintillation Using Fixed SBAS Links at Low Latitudes}
  {Scientific Instrumentation and Challenges}

% row 47
\oral
  {Eliah F. M. T. São Sabbas}
  {Transient Luminous Event (TLE) Observations with the LEONA Network}
  {Scientific Instrumentation and Challenges}

% row 48
\ouvinte{Adebayo Samuel Sunday}

% row 49
\ouvinte{Adriano P. Almeida}

% row 50
\ouvinte{Adriany Rodrigues Barbosa}

% row 51
\ouvinte{Agri Faturahman}

% row 52
\ouvinte{Ângela Santos}

% row 53
\ouvinte{Daiene de Mello Schaefer}

% row 54
\ouvinte{Dave DELVA}

% row 55
\ouvinte{Deysi Veronica Cornejo Espinoza}

% row 56
\ouvinte{Edith Tueros Cuadros}

% row 57
\ouvinte{Elias de Vargas Becker}

% row 58
\ouvinte{Emilly Silva Rosa}

% row 59
\ouvinte{Gbadebo Seun Adeniran}

% row 60
\ouvinte{Hadassa Raquel Peixoto Jácome}

% row 61
\ouvinte{Ian Gabriel Oliveira da Silva}

% row 62
\ouvinte{Igor Gabriel Gonçalves Vicente}

% row 63
\ouvinte{Inez Staciarini Batista}

% row 64
\ouvinte{Isela Leonor Vasquez Panduro}

% row 65
\ouvinte{Jaziel Felipe Braga Campelo}

% row 66
\ouvinte{Jose Paulo Marchezi}

% row 67
\ouvinte{Luan Mateus Bocalan Vogás}

% row 68
\ouvinte{Peter Taiwo Muka}

% row 69
\ouvinte{Renan Yule Queiroz Ferreira}

% row 70
\ouvinte{Rochelin YAYO}

% row 71
\ouvinte{Simbiat Oamen}

% row 72
\ouvinte{Veera Kumar}

% row 73
\ouvinte{vinicius deggeroni}

\end{document}
